\section{Development of the Ordinary: a dated chronology}

This table gathers the datable acts and witnesses on which the exposition
draws, so that a reader can see at once which claims rest on a document, which
on a reported reconstruction, and which on the collated facsimile itself. It is
an orientation aid, not an argument; every entry is discussed where it belongs
in the exposition, and the character of the evidence is stated in the third
column rather than smoothed away.

\begin{chronotable}
\unitrow{c.~390}{\work{De sacramentis} IV, 5--6, preached at Milan}{Latin
eucharistic prayer with \latincue{ascriptam, ratam, rationabilem,
acceptabilem}; \latincue{Qui pridie quam pateretur \dots\ benedixit, fregit};
\latincue{memores} with passion, resurrection and ascension; the offering
carried \latincue{per manus angelorum tuorum} to the altar on high; Abel,
Abraham, Melchisedech in that order. Verbally different from the Roman text at
every point of comparison. Read at \work{PL} 16, coll.~462--464.}
\unitrow{590--604}{Pontificate of Gregory the Great}{Fortescue names three
changes as his --- the litany's petitions dropped from the daily
\latincue{Kyrie}, an addition to \latincue{Hanc igitur}, and the transfer of
the Lord's Prayer to its place before the Communion --- and records the
constant tradition, with Benedict XIV's dictum, that no pope changed the Canon
after him. Reported reconstruction and reported tradition.}
\unitrow{687--701}{Pontificate of Sergius I}{The \work{Liber Pontificalis},
as Fortescue reports it, has Sergius order the \latincue{Agnus Dei} sung by
clergy and people at the breaking of the Lord's body. Reported.}
\unitrow{1014}{Coronation of Henry II at Rome, 14 February}{Berno of Reichenau
reports that no creed was sung and that Benedict VIII afterwards ordered it
sung after the Gospel at Rome. Fortescue records that most authors accept the
account and names in a footnote those who do not. Probable, not secure.}
\unitrow{Twelfth century}{John Beleth}{Describes the threefold
\latincue{Agnus Dei} ending \latincue{dona nobis pacem}; Innocent III notes
churches, the Lateran among them, which kept \latincue{miserere nobis} three
times --- the shape the Requiem form still has. Reported.}
\unitrow{Fourteenth century}{The fourteenth Roman \work{Ordo}}{Fortescue
identifies its offertory prayers as the set the revisers of 1570 adopted, ``all
taken or adapted from various, mostly non-Roman sources'', and traces
\latincue{Suscipe sancte Pater} to the prayerbook of Charles the Bald.
The manuscript is defensibly dated 846--869; Fortescue's 875--877 are
Charles's imperial years. His broader reconstruction is reported rather than
independently established here.}
\unitrow{1502}{Burchard's ceremonial}{Allows the last Gospel, until then part
of the priest's thanksgiving on the way to the sacristy, to be said at the
altar. Reported.}
\unitrow{17 September 1562}{Council of Trent, session XXII}{Chapter III on
Masses in honour of the saints, whose closing clause is the closing clause of
\latincue{Suscipe, sancta Trinitas}; chapter IV \latincue{De canone missae},
declaring the Canon free from error and composed of the Lord's words, apostolic
traditions and papal institutions; chapter V on the low and louder voice;
chapter VII on the water in the chalice; canon VI against the assertion that
the Canon contains errors; canon IX against the condemnation of the low-voiced
Canon. Read at the page images of the Tauchnitz Latin and Waterworth's English.}
\unitrow{14 July 1570}{\latincue{Quo primum}, Pius V}{Imposes the restored
Missal, exempting rites approved from their first institution or continuously
observed for more than two hundred years. Read at the facsimile's own front
matter. This is the edition in which the last Gospel first enters the Mass and
in which the medieval offertory and Communion prayers are fixed.}
\unitrow{7 July 1604}{\latincue{Cum sanctissimum}, Clement VIII}{Records that
Pius V had forbidden addition or subtraction under severe penalties, that
corruptions had nevertheless entered printed Missals, and restores the book.
Read at the facsimile's front matter.}
\unitrow{2 September 1634}{\latincue{Si quid est in rebus humanis}, Urban
VIII}{Restores rubrics that had degenerated from ancient use and explains those
less clear to readers; the whole collated with the Vulgate. Read at the
facsimile's front matter.}
\unitrow{25 July 1960}{\latincue{Rubricarum instructum}, John XXIII}{Approves
the new code of rubrics. The 1962 Missal's own promulgating decree gives the
date as 23 July; the Congregation's general decree, printed later in the same
volume, gives 25 July, which alone is consistent with that decree's
\latincue{postero die}. The discrepancy is internal to the book and is not
resolved here.}
\unitrow{26 July 1960}{General decree of the Sacred Congregation of
Rites}{Promulgates the code, orders it inserted in new editions of the Breviary
and Missal, and requires observance from 1 January 1961. This is the source of
nearly every rubric quoted in this study, including nn.~424--425, 500--503 and
511--512.}
\unitrow{23 June 1962}{Decree declaring this Vatican edition typical}{States
that a new edition of the Missal fully accommodated to the code of rubrics
could scarcely fail to be prepared. Read at the page image of the facsimile's
front matter.}
\unitrow{13 November 1962}{\latincue{De S. Ioseph nomine Canoni Missae
inserendo}}{The Sacred Congregation of Rites decrees that after
\latincue{Communicantes \dots\ Domini nostri Iesu Christi} the words
\latincue{sed et beati Ioseph eiusdem Virginis Sponsi} be added, extending the
prescription to the days with a proper \latincue{Communicantes}. John XXIII's
decision was announced to the Council fathers in the Vatican basilica the same
day; the decree orders it in force from 8 December 1962. \work{AAS} 54 (1962),
p.~873.}
\unitrow{8 December 1962}{The insertion takes effect}{The facsimile collated
for this study prints the clause both in the ordinary \latincue{Communicantes}
at n.~1090 and in the shared continuation of the five proper forms at printed
p.~305. The 1861 hand missal, printed a century earlier, does not.}
\end{chronotable}
